\section{Teorema de Bayes}
\label{cap:teoremaBayes}

O Teorema de Bayes, resultado dos estudos matemáticos de Thomas Bayes e Pierre-Simon Laplace, é um teorema que permite atualizar a probabilidade de um evento a partir de novas informações disponíveis. Além disso, ele possibilita calcular a probabilidade de uma hipótese ser verdadeira sabendo que determinado evento ocorreu, mesmo quando se conhece apenas a probabilidade de esse evento ocorrer caso a hipótese seja verdadeira. Por essa razão, o teorema é amplamente utilizado em situações que envolvem tomada de decisão sob incerteza e interpretação de evidências.

Matematicamente, o Teorema de Bayes pode ser expresso pela fórmula:

\begin{equation}
    P(A|B) = \frac{P(B|A) \cdot P(A)}{P(B)}
    \label{eq:teorema_bayes}
\end{equation}

%TODO - este texto está muito chucro. Melhorar!

em que $P(A)$ representa a probabilidade inicial da hipótese, $P(B|A)$ corresponde à probabilidade de observar a evidência caso a hipótese seja verdadeira e $P(B)$ representa a probabilidade total da evidência ocorrer. O resultado $P(A|B)$ é chamado de probabilidade posterior, pois corresponde à nova estimativa da hipótese após a incorporação das informações observadas. Esse processo de atualização é extremamente importante para cálculos estatísticos, sendo descrito por Harold Jeffreys como o equivalente para a probabilidade do teorema de Pitágoras para a geometria.

O Teorema de Bayes possui diversas aplicações práticas em áreas como medicina, inteligência artificial, aprendizado de máquina, sistemas de recomendação e detecção de fraudes. Em exames médicos, por exemplo, um teste positivo não significa necessariamente que um paciente possui determinada doença, sendo necessário considerar também a frequência da doença na população. Assim, o teorema permite combinar as informações iniciais com os resultados obtidos para produzir estimativas mais precisas sobre a ocorrência de um evento.

% TODO - melhorar este código, gerando randomicamente as coisas com Python e aperfeiçoar a exemplificação do problema de Bayes
\begin{lstlisting}[language=Python]

A = float(input("Qual a probablidade de A: \n"))
B = float(input("Qual a probablidade de B: \n"))
AB = float(input("Qual a probablidade de A|B: \n"))

print("A possibilidade de B dado A é: ", AB * B / A, "")

\end{lstlisting}

